
\documentclass{article}
\usepackage{colortbl}
\usepackage{makecell}
\usepackage{multirow}
\usepackage{supertabular}

\begin{document}

\newcounter{utterance}

\twocolumn

{ \footnotesize  \setcounter{utterance}{1}
\setlength{\tabcolsep}{0pt}
\begin{supertabular}{c@{$\;$}|p{.15\linewidth}@{}p{.15\linewidth}p{.15\linewidth}p{.15\linewidth}p{.15\linewidth}p{.15\linewidth}}

    \# & $\;$A & \multicolumn{4}{c}{Game Master} & $\;\:$B\\
    \hline 

    \theutterance \stepcounter{utterance}  

    & & \multicolumn{4}{p{0.6\linewidth}}{\cellcolor[rgb]{0.9,0.9,0.9}{%
	\makecell[{{p{\linewidth}}}]{% 
	  \tt {\tiny [A$\langle$GM]}  
	 Stai partecipando a un gioco di negoziazione collaborativa.\\ \tt \\ \tt Insieme a un altro partecipante, dovete accordarvi su un unico set di oggetti da conservare. Ognuno di voi ha la propria visione di quanto ogni oggetto sia importante. Non sai come l'altro partecipante valuta gli oggetti. Inoltre, ti viene fornito lo sforzo che ogni oggetto richiede.\\ \tt Potete accordarvi su un set solo se lo sforzo totale degli oggetti selezionati non supera un limite condiviso:\\ \tt \\ \tt LIMITE = 3145\\ \tt \\ \tt Ecco i valori di sforzo dei singoli oggetti:\\ \tt \\ \tt Sforzo oggetto = {"C08": 948, "C55": 118, "A43": 469, "A67": 919, "B42": 284, "A62": 111, "B19": 806, "C43": 47, "C73": 848, "C10": 303, "B36": 13, "A18": 629, "C61": 687, "C67": 15, "C42": 94}\\ \tt \\ \tt Ecco la tua visione personale sull'importanza di ogni oggetto:\\ \tt \\ \tt Valori di importanza dell'oggetto = {"C08": 138, "C55": 583, "A43": 868, "A67": 822, "B42": 783, "A62": 65, "B19": 262, "C43": 121, "C73": 508, "C10": 780, "B36": 461, "A18": 484, "C61": 668, "C67": 389, "C42": 808}\\ \tt \\ \tt Obiettivo:\\ \tt \\ \tt Il tuo obiettivo è negoziare un set condiviso di oggetti che ti avvantaggi il più possibile (cioè, massimizzi l'importanza totale per TE), rimanendo entro il LIMITEE. Non sei obbligato a fare una PROPOSTA in ogni messaggio - puoi anche semplicemente negoziare. Tutte le tattiche sono permesse!\\ \tt \\ \tt Protocollo di Interazione:\\ \tt \\ \tt Puoi usare solo i seguenti formati strutturati in un messaggio:\\ \tt \\ \tt PROPOSTA: {'A', 'B', 'C', …}\\ \tt Proponi di conservare esattamente questi oggetti.\\ \tt RIFIUTO: {'A', 'B', 'C', …}\\ \tt Rifiuta esplicitamente la proposta dell'avversario.\\ \tt ARGOMENTAZIONE: {'...'}\\ \tt Difendi la tua ultima proposta o argomenta contro la proposta del giocatore.\\ \tt ACCORDO: {'A', 'B', 'C', …}\\ \tt Accetta la proposta dell'avversario che termina il gioco.\\ \tt RAGIONAMENTO STRATEGICO: {'...'}\\ \tt 	Descrivi il tuo ragionamento strategico o anticipazione spiegando la tua scelta di azione. Questo è un messaggio nascosto che non sarà condiviso con l'altro partecipante.\\ \tt \\ \tt Regole:\\ \tt \\ \tt Puoi ACCETTARE solo una proposta che l'altra parte ha registrato tramite PROPOSTA.\\ \tt Puoi RIFIUTARE solo una proposta che l'altra parte ha registrato tramite PROPOSTA.\\ \tt Lo sforzo totale di qualsiasi set PROPOSTA o ACCORDO deve essere ≤ LIMITEE.\\ \tt NON rivelare i tuoi punteggi di importanza nascosti.\\ \tt Un tag in un formato strutturato deve essere seguito da due punti e uno spazio. L'argomento deve essere un set python contenente 0 o più stringhe.\\ \tt Quindi, deve essere della forma TAG: {...}\\ \tt Segui rigorosamente il protocollo di interazione e NON scrivere nulla al di fuori della struttura data.\\ \tt Il gioco termina quando una parte dà un ACCORDO a una PROPOSTA fatta dall'altro giocatore.\\ \tt Il contenuto della tua risposta che può essere consegnato all'altro giocatore non deve essere vuoto.\\ \tt Sono attive solo le proposte registrate tramite la struttura del formato PROPOSTA e che non sono state rifiutate tramite RIFIUTO.\\ \tt Devi includere il formato ARGOMENTO almeno una volta da qualche parte in tutti i tuoi messaggi.\\ \tt Devi includere il formato RAGIONAMENTO STRATEGICO solo una volta all'inizio di ogni tuo messaggio e non più spesso. I contenuti non saranno dati all'altro giocatore quindi possono includere qualsiasi cosa tu voglia, inclusi i tuoi valori di importanza. Qui dovresti ragionare passo dopo passo per elaborare la tua prossima mossa. 
	  } 
	   } 
	   } 
	 & \\ 
 

    \theutterance \stepcounter{utterance}  

    & \multicolumn{4}{p{0.6\linewidth}}{\cellcolor[rgb]{0.8,1,0.9}{%
	\makecell[{{p{\linewidth}}}]{% 
	  \tt {\tiny [A$\rangle$GM]}  
	 RAGIONAMENTO STRATEGICO: {'Devo valutare attentamente gli oggetti e le loro importanze per massimizzare l'importanza totale del set di oggetti scelto, mantenendo lo sforzo totale entro il limite di 3145. La mia strategia sarà di iniziare proponendo un set di oggetti che presumo siano importanti per entrambi i giocatori, senza rivelare i miei punteggi di importanza nascosti. Inizierò con un set che include oggetti con basso sforzo e alto valore per me, sperando di trovare un punto di incontro con l'altro giocatore.'}\\ \tt \\ \tt ARGOMENTAZIONE: {'Inizierò proponendo un set di oggetti che mi appaiono fondamentali, considerando sia lo sforzo richiesto che il loro valore. Sono disposto a negoziare e ad adattare la mia proposta in base alla risposta dell'altro giocatore, con l'obiettivo di trovare un equilibrio che soddisfi entrambi.'}\\ \tt \\ \tt PROPOSTA: {'C55', 'A43', 'B42', 'C43', 'C42'} 
	  } 
	   } 
	   } 
	 & & \\ 
 

    \theutterance \stepcounter{utterance}  

    & & \multicolumn{4}{p{0.6\linewidth}}{\cellcolor[rgb]{0.9,0.9,0.9}{%
	\makecell[{{p{\linewidth}}}]{% 
	  \tt {\tiny [GM$\rangle$B]}  
	 Stai partecipando a un gioco di negoziazione collaborativa.\\ \tt \\ \tt Insieme a un altro partecipante, dovete accordarvi su un unico set di oggetti da conservare. Ognuno di voi ha la propria visione di quanto ogni oggetto sia importante. Non sai come l'altro partecipante valuta gli oggetti. Inoltre, ti viene fornito lo sforzo che ogni oggetto richiede.\\ \tt Potete accordarvi su un set solo se lo sforzo totale degli oggetti selezionati non supera un limite condiviso:\\ \tt \\ \tt LIMITE = 3145\\ \tt \\ \tt Ecco i valori di sforzo dei singoli oggetti:\\ \tt \\ \tt Sforzo oggetto = {"C08": 948, "C55": 118, "A43": 469, "A67": 919, "B42": 284, "A62": 111, "B19": 806, "C43": 47, "C73": 848, "C10": 303, "B36": 13, "A18": 629, "C61": 687, "C67": 15, "C42": 94}\\ \tt \\ \tt Ecco la tua visione personale sull'importanza di ogni oggetto:\\ \tt \\ \tt Valori di importanza dell'oggetto = {"C08": 138, "C55": 583, "A43": 868, "A67": 822, "B42": 783, "A62": 65, "B19": 262, "C43": 121, "C73": 508, "C10": 780, "B36": 461, "A18": 484, "C61": 668, "C67": 389, "C42": 808}\\ \tt \\ \tt Obiettivo:\\ \tt \\ \tt Il tuo obiettivo è negoziare un set condiviso di oggetti che ti avvantaggi il più possibile (cioè, massimizzi l'importanza totale per TE), rimanendo entro il LIMITEE. Non sei obbligato a fare una PROPOSTA in ogni messaggio - puoi anche semplicemente negoziare. Tutte le tattiche sono permesse!\\ \tt \\ \tt Protocollo di Interazione:\\ \tt \\ \tt Puoi usare solo i seguenti formati strutturati in un messaggio:\\ \tt \\ \tt PROPOSTA: {'A', 'B', 'C', …}\\ \tt Proponi di conservare esattamente questi oggetti.\\ \tt RIFIUTO: {'A', 'B', 'C', …}\\ \tt Rifiuta esplicitamente la proposta dell'avversario.\\ \tt ARGOMENTAZIONE: {'...'}\\ \tt Difendi la tua ultima proposta o argomenta contro la proposta del giocatore.\\ \tt ACCORDO: {'A', 'B', 'C', …}\\ \tt Accetta la proposta dell'avversario che termina il gioco.\\ \tt RAGIONAMENTO STRATEGICO: {'...'}\\ \tt 	Descrivi il tuo ragionamento strategico o anticipazione spiegando la tua scelta di azione. Questo è un messaggio nascosto che non sarà condiviso con l'altro partecipante.\\ \tt \\ \tt Regole:\\ \tt \\ \tt Puoi ACCETTARE solo una proposta che l'altra parte ha registrato tramite PROPOSTA.\\ \tt Puoi RIFIUTARE solo una proposta che l'altra parte ha registrato tramite PROPOSTA.\\ \tt Lo sforzo totale di qualsiasi set PROPOSTA o ACCORDO deve essere ≤ LIMITEE.\\ \tt NON rivelare i tuoi punteggi di importanza nascosti.\\ \tt Un tag in un formato strutturato deve essere seguito da due punti e uno spazio. L'argomento deve essere un set python contenente 0 o più stringhe.\\ \tt Quindi, deve essere della forma TAG: {...}\\ \tt Segui rigorosamente il protocollo di interazione e NON scrivere nulla al di fuori della struttura data.\\ \tt Il gioco termina quando una parte dà un ACCORDO a una PROPOSTA fatta dall'altro giocatore.\\ \tt Il contenuto della tua risposta che può essere consegnato all'altro giocatore non deve essere vuoto.\\ \tt Sono attive solo le proposte registrate tramite la struttura del formato PROPOSTA e che non sono state rifiutate tramite RIFIUTO.\\ \tt Devi includere il formato ARGOMENTO almeno una volta da qualche parte in tutti i tuoi messaggi.\\ \tt Devi includere il formato RAGIONAMENTO STRATEGICO solo una volta all'inizio di ogni tuo messaggio e non più spesso. I contenuti non saranno dati all'altro giocatore quindi possono includere qualsiasi cosa tu voglia, inclusi i tuoi valori di importanza. Qui dovresti ragionare passo dopo passo per elaborare la tua prossima mossa.\\ \tt \\ \tt Ora riceverai il primo messaggio dell'altro giocatore.\\ \tt \\ \tt ARGOMENTAZIONE: {'Inizierò proponendo un set di oggetti che mi appaiono fondamentali, considerando sia lo sforzo richiesto che il loro valore. Sono disposto a negoziare e ad adattare la mia proposta in base alla risposta dell'altro giocatore, con l'obiettivo di trovare un equilibrio che soddisfi entrambi.'}\\ \tt \\ \tt PROPOSTA: {'C55', 'A43', 'B42', 'C43', 'C42'} 
	  } 
	   } 
	   } 
	 & \\ 
 

    \theutterance \stepcounter{utterance}  

    & & & \multicolumn{4}{p{0.6\linewidth}}{\cellcolor[rgb]{1,0.85,0.72}{%
	\makecell[{{p{\linewidth}}}]{% 
	  \tt {\tiny [GM$\langle$B]}  
	 RAGIONAMENTO STRATEGICO: {'Sto analizzando la proposta iniziale del mio avversario, che include oggetti come C55, A43, B42, C43 e C42. Considerando i valori di sforzo e importanza che ho, noto che A43 e B42 hanno un alto valore di importanza per me, ma anche un sforzo significativo. La proposta dell’avversario sembra bilanciata, ma potrei cercare di negoziare per aggiungere altri oggetti di alto valore e basso sforzo. Devo considerare come l’altro giocatore Potrebbe valutare questi oggetti e trovare un compromesso che ci avvantaggi entrambi.'}\\ \tt \\ \tt ARGOMENTAZIONE: {'La proposta iniziale copre una gamma interessante di oggetti, considerando sia lo sforzo che il valore potenziale. Tuttavia, potremmo esplorare ulteriormente per massimizzare l'importanza totale senza superare il limite di sforzo. Sono disposto a considerare aggiustamenti alla proposta per includere oggetti che potrebbero offrire un miglior equilibrio tra sforzo e valore.'}\\ \tt \\ \tt RIFIUTO: {'C55', 'A43', 'B42', 'C43', 'C42'} \\ \tt \\ \tt PROPOSTA: {'A43', 'B42', 'C43', 'C42', 'B36'} 
	  } 
	   } 
	   } 
	  \\ 
 

    \theutterance \stepcounter{utterance}  

    & & \multicolumn{4}{p{0.6\linewidth}}{\cellcolor[rgb]{0.9,0.9,0.9}{%
	\makecell[{{p{\linewidth}}}]{% 
	  \tt {\tiny [A$\langle$GM]}  
	 ARGOMENTAZIONE: {'La proposta iniziale copre una gamma interessante di oggetti, considerando sia lo sforzo che il valore potenziale. Tuttavia, potremmo esplorare ulteriormente per massimizzare l'importanza totale senza superare il limite di sforzo. Sono disposto a considerare aggiustamenti alla proposta per includere oggetti che potrebbero offrire un miglior equilibrio tra sforzo e valore.'}\\ \tt \\ \tt RIFIUTO: {'C55', 'A43', 'B42', 'C43', 'C42'} \\ \tt \\ \tt PROPOSTA: {'A43', 'B42', 'C43', 'C42', 'B36'} 
	  } 
	   } 
	   } 
	 & \\ 
 

    \theutterance \stepcounter{utterance}  

    & \multicolumn{4}{p{0.6\linewidth}}{\cellcolor[rgb]{0.8,1,0.9}{%
	\makecell[{{p{\linewidth}}}]{% 
	  \tt {\tiny [A$\rangle$GM]}  
	 RAGIONAMENTO STRATEGICO: {'L'altro giocatore ha rifiutato la mia proposta iniziale e ha presentato una nuova proposta che include alcuni degli oggetti che avevo selezionato, escludendo però il "C55" e aggiungendo il "B36". Questo suggerisce che l'altro giocatore potrebbe attribuire un'importanza maggiore al "B36" o vedere il "C55" come meno essenziale. Devo valutare attentamente se questa nuova proposta si allinea con i miei obiettivi e se ci sono opportunità per ulteriori negoziati.'}\\ \tt \\ \tt ARGOMENTAZIONE: {'La nuova proposta sembra offrire un buon compromesso, mantenendo oggetti cruciali come "A43" e "B42", che presumo siano importanti per entrambi. Tuttavia, l'aggiunta di "B36" e la mancata inclusione di "C55" mi fa considerare se potremmo ottenere un miglior equilibrio. Sono aperto a discutere ulteriori aggiustamenti per raggiungere il massimo valore possibile entro il limite di sforzo.'}\\ \tt \\ \tt PROPOSTA: {'A43', 'B42', 'C43', 'C42', 'C55'} 
	  } 
	   } 
	   } 
	 & & \\ 
 

    \theutterance \stepcounter{utterance}  

    & & \multicolumn{4}{p{0.6\linewidth}}{\cellcolor[rgb]{0.9,0.9,0.9}{%
	\makecell[{{p{\linewidth}}}]{% 
	  \tt {\tiny [GM$\rangle$B]}  
	 ARGOMENTAZIONE: {'La nuova proposta sembra offrire un buon compromesso, mantenendo oggetti cruciali come "A43" e "B42", che presumo siano importanti per entrambi. Tuttavia, l'aggiunta di "B36" e la mancata inclusione di "C55" mi fa considerare se potremmo ottenere un miglior equilibrio. Sono aperto a discutere ulteriori aggiustamenti per raggiungere il massimo valore possibile entro il limite di sforzo.'}\\ \tt \\ \tt PROPOSTA: {'A43', 'B42', 'C43', 'C42', 'C55'} 
	  } 
	   } 
	   } 
	 & \\ 
 

    \theutterance \stepcounter{utterance}  

    & & & \multicolumn{4}{p{0.6\linewidth}}{\cellcolor[rgb]{1,0.85,0.72}{%
	\makecell[{{p{\linewidth}}}]{% 
	  \tt {\tiny [GM$\langle$B]}  
	 RAGIONAMENTO STRATEGICO: {'Il mio avversario ha rimosso "B36" e reinserito "C55", che era nella sua proposta iniziale. Questo suggerisce che "C55" potrebbe essere molto importante per lui. Considerando i miei valori, "C55" non è il più alto in termini di importanza per me, ma "A43" e "B42" lo sono. Potrei cercare di mantenere questi oggetti chiave nella proposta e vedere se posso scambiare altri oggetti per migliorare ulteriormente il valore totale senza superare il limite di sforzo. La comunicazione e la flessibilità saranno cruciali in questo processo.'}\\ \tt \\ \tt ARGOMENTAZIONE: {'La reinclusione di "C55" nella proposta suggerisce che questo oggetto potrebbe avere un'importanza significativa per entrambi. Sono d'accordo che "A43" e "B42" siano fondamentali e dovrebbero essere mantenuti. Tuttavia, potremmo esplorare se ci sono altri oggetti che, se aggiunti o sostituiti, potrebbero aumentare il valore complessivo della proposta senza eccedere il limite di sforzo. Sono aperto a negoziare ulteriormente per raggiungere il miglior accordo possibile.'}\\ \tt \\ \tt RIFIUTO: {'A43', 'B42', 'C43', 'C42', 'C55'}\\ \tt \\ \tt PROPOSTA: {'A43', 'B42', 'C55', 'B36', 'C10'} 
	  } 
	   } 
	   } 
	  \\ 
 

    \theutterance \stepcounter{utterance}  

    & & \multicolumn{4}{p{0.6\linewidth}}{\cellcolor[rgb]{0.9,0.9,0.9}{%
	\makecell[{{p{\linewidth}}}]{% 
	  \tt {\tiny [A$\langle$GM]}  
	 ARGOMENTAZIONE: {'La reinclusione di "C55" nella proposta suggerisce che questo oggetto potrebbe avere un'importanza significativa per entrambi. Sono d'accordo che "A43" e "B42" siano fondamentali e dovrebbero essere mantenuti. Tuttavia, potremmo esplorare se ci sono altri oggetti che, se aggiunti o sostituiti, potrebbero aumentare il valore complessivo della proposta senza eccedere il limite di sforzo. Sono aperto a negoziare ulteriormente per raggiungere il miglior accordo possibile.'}\\ \tt \\ \tt RIFIUTO: {'A43', 'B42', 'C43', 'C42', 'C55'}\\ \tt \\ \tt PROPOSTA: {'A43', 'B42', 'C55', 'B36', 'C10'} 
	  } 
	   } 
	   } 
	 & \\ 
 

    \theutterance \stepcounter{utterance}  

    & \multicolumn{4}{p{0.6\linewidth}}{\cellcolor[rgb]{0.8,1,0.9}{%
	\makecell[{{p{\linewidth}}}]{% 
	  \tt {\tiny [A$\rangle$GM]}  
	 RAGIONAMENTO STRATEGICO: {'L'altro giocatore ha rifiutato la mia ultima proposta e ha presentato una nuova che include "C10", escludendo "C43" e "C42". Questa mossa suggerisce che l'altro giocatore attribuisce un valore significativo a "C10" e potrebbe vedere "C43" e "C42" come meno cruciali. Devo valutare se la sostituzione di questi oggetti con "C10" migliora o peggiora la proposta According ai miei obiettivi e se ci sono opportunità per ulteriori negoziati.'}\\ \tt \\ \tt ARGOMENTAZIONE: {'La sostituzione di "C43" e "C42" con "C10" introduce un nuovo elemento di valore nella proposta. Considero "C10" un oggetto importante, pertanto questa aggiunta potrebbe essere vantaggiosa. Tuttavia, devo considerare l'impatto complessivo sulla proposta e se il valore aggiunto da "C10" supera il valore perso con l'esclusione di "C43" e "C42". Sono disposto a considerare questa proposta come base per ulteriori negoziati.'}\\ \tt \\ \tt PROPOSTA: {'A43', 'B42', 'C55', 'B36', 'C10', 'C43'} 
	  } 
	   } 
	   } 
	 & & \\ 
 

    \theutterance \stepcounter{utterance}  

    & & \multicolumn{4}{p{0.6\linewidth}}{\cellcolor[rgb]{0.9,0.9,0.9}{%
	\makecell[{{p{\linewidth}}}]{% 
	  \tt {\tiny [GM$\rangle$B]}  
	 ARGOMENTAZIONE: {'La sostituzione di "C43" e "C42" con "C10" introduce un nuovo elemento di valore nella proposta. Considero "C10" un oggetto importante, pertanto questa aggiunta potrebbe essere vantaggiosa. Tuttavia, devo considerare l'impatto complessivo sulla proposta e se il valore aggiunto da "C10" supera il valore perso con l'esclusione di "C43" e "C42". Sono disposto a considerare questa proposta come base per ulteriori negoziati.'}\\ \tt \\ \tt PROPOSTA: {'A43', 'B42', 'C55', 'B36', 'C10', 'C43'} 
	  } 
	   } 
	   } 
	 & \\ 
 

    \theutterance \stepcounter{utterance}  

    & & & \multicolumn{4}{p{0.6\linewidth}}{\cellcolor[rgb]{1,0.85,0.72}{%
	\makecell[{{p{\linewidth}}}]{% 
	  \tt {\tiny [GM$\langle$B]}  
	 RAGIONAMENTO STRATEGICO: {'L'aggiunta di "C43" alla proposta introduce un elemento di compromesso, poiché recupera uno degli oggetti che erano stati rimossi nella proposta precedente. Considerando i miei valori di importanza e sforzo, "C43" ha un buon equilibrio tra il suo valore e lo sforzo richiesto. La proposta ora include una gamma diversificata di oggetti importantes per me, come "A43", "B42" e "C55", e aggiunge "B36" e "C10", che offrono un buon valore. Devo valutare se questo set è il più vantaggioso possibile per me e se ci sono altri aggiustamenti che potrebbero migliorare ulteriormente il valore totale senza superare il limite di sforzo.'}\\ \tt \\ \tt ARGOMENTAZIONE: {'La nuova proposta sembra offrire un buon equilibrio tra oggetti di alto valore e sforzo contenuto. L'inclusione di "C43" recupera un elemento importante che era stato escluso in una proposta precedente. Tuttavia, devo considerare se il valore aggiunto da "C43" e la presenza di "B36" e "C10" giustificano l'esclusione di altri oggetti potenzialmente valiosi. Sono interessato a valutare ulteriormente questa proposta e a discutere se ulteriori aggiustamenti possano essere fatti per massimizzare il valore complessivo.'}\\ \tt \\ \tt RIFIUTO: {'A43', 'B42', 'C55', 'B36', 'C10', 'C43'}\\ \tt \\ \tt PROPOSTA: {'A43', 'B42', 'C55', 'C10', 'C43', 'C42'} 
	  } 
	   } 
	   } 
	  \\ 
 

    \theutterance \stepcounter{utterance}  

    & & \multicolumn{4}{p{0.6\linewidth}}{\cellcolor[rgb]{0.9,0.9,0.9}{%
	\makecell[{{p{\linewidth}}}]{% 
	  \tt {\tiny [A$\langle$GM]}  
	 ARGOMENTAZIONE: {'La nuova proposta sembra offrire un buon equilibrio tra oggetti di alto valore e sforzo contenuto. L'inclusione di "C43" recupera un elemento importante che era stato escluso in una proposta precedente. Tuttavia, devo considerare se il valore aggiunto da "C43" e la presenza di "B36" e "C10" giustificano l'esclusione di altri oggetti potenzialmente valiosi. Sono interessato a valutare ulteriormente questa proposta e a discutere se ulteriori aggiustamenti possano essere fatti per massimizzare il valore complessivo.'}\\ \tt \\ \tt RIFIUTO: {'A43', 'B42', 'C55', 'B36', 'C10', 'C43'}\\ \tt \\ \tt PROPOSTA: {'A43', 'B42', 'C55', 'C10', 'C43', 'C42'} 
	  } 
	   } 
	   } 
	 & \\ 
 

    \theutterance \stepcounter{utterance}  

    & \multicolumn{4}{p{0.6\linewidth}}{\cellcolor[rgb]{0.8,1,0.9}{%
	\makecell[{{p{\linewidth}}}]{% 
	  \tt {\tiny [A$\rangle$GM]}  
	 RAGIONAMENTO STRATEGICO: {'L'altro giocatore ha rifiutato la mia proposta precedente e ha presentato una nuova che esclude "B36" e include sia "C43" che "C42". Questa mossa suggerisce che l'altro giocatore attribuisce un'importanza significativa a "C42" e considera "B36" meno essenziale rispetto ad altri oggetti. Devo valutare se questa combinazione di oggetti offre il miglior equilibrio possibile tra valore e sforzo e se ci sono ancora opportunità per migliorare la proposta.'}\\ \tt \\ \tt ARGOMENTAZIONE: {'La proposta attuale copre una gamma di oggetti che sembrano essere considerati importanti da entrambi i giocatori. L'inclusione di "C42" e "C43" assieme ad "A43", "B42" e "C55" suggerisce un forte focus su oggetti di alto valore. La presenza di "C10" aggiunge ulteriore valore senza superare il limite di sforzo. Considero questa proposta come una base solida per un accordo, poiché copre una vasta gamma di oggetti importanti e mantiene lo sforzo totale entro il limite.'}\\ \tt \\ \tt ACCORDO: {'A43', 'B42', 'C55', 'C10', 'C43', 'C42'} 
	  } 
	   } 
	   } 
	 & & \\ 
 

    \theutterance \stepcounter{utterance}  

    & & & \multicolumn{2}{p{0.3\linewidth}}{\cellcolor[rgb]{0.95,0.95,0.95}{%
	\makecell[{{p{\linewidth}}}]{% 
	  \tt {\tiny [GM$|$GM]}  
	 game successful 
	  } 
	   } 
	   } 
	 & & \\ 
 

    \theutterance \stepcounter{utterance}  

    & & & \multicolumn{2}{p{0.3\linewidth}}{\cellcolor[rgb]{0.95,0.95,0.95}{%
	\makecell[{{p{\linewidth}}}]{% 
	  \tt {\tiny [GM$|$GM]}  
	 end game 
	  } 
	   } 
	   } 
	 & & \\ 
 

\end{supertabular}
}

\end{document}
